A partir de los grafos Mapper reconciliados, la interpretación principal del reporte se concentró en seis grafos con lente \texttt{pca2}. Esta selección permite comparar espacios físicos, orbitales, térmicos y conjuntos bajo una misma configuración de cubierta, \texttt{cubes10\_overlap0p35}. Los grafos seleccionados no cumplen todos la misma función: algunos se usan como resultados primarios, otros como controles metodológicos y uno como caso cauteloso por su alta dependencia de imputación.

Los grafos primarios son aquellos que permiten interpretar directamente una dimensión central del problema: la estructura física mínima, la estructura orbital y la estructura conjunta sin densidad. Los grafos de control permiten evaluar el efecto de agregar \texttt{pl\_dens}, una variable físicamente útil pero algebraicamente dependiente de masa y radio en gran parte del catálogo. El grafo térmico se conserva como resultado de interés, pero se etiqueta como \textit{cautionary} porque su estructura depende fuertemente de variables imputadas.

\begin{table}[H]
\centering
\small
\begin{tabular}{p{0.48\textwidth}p{0.18\textwidth}p{0.24\textwidth}}
\toprule
\textbf{Grafo seleccionado} & \textbf{Rol} & \textbf{Lectura principal} \\
\midrule
\texttt{phys\_min\_pca2\_cubes10\_overlap0p35} & Primario & Estructura física mínima usando radio y masa. \\
\texttt{phys\_density\_pca2\_cubes10\_overlap0p35} & Control & Efecto de agregar densidad al espacio físico. \\
\texttt{orbital\_pca2\_cubes10\_overlap0p35} & Primario & Estructura orbital principal. \\
\texttt{joint\_no\_density\_pca2\_cubes10\_overlap0p35} & Primario & Estructura conjunta sin densidad derivada. \\
\texttt{joint\_pca2\_cubes10\_overlap0p35} & Control & Efecto de agregar densidad al espacio conjunto. \\
\texttt{thermal\_pca2\_cubes10\_overlap0p35} & Cautionary & Estructura térmica con alta dependencia de imputación. \\
\bottomrule
\end{tabular}
\caption{Grafos \texttt{pca2} seleccionados para la interpretación principal del reporte.}
\label{tab:selected_pca2_graphs}
\end{table}

Entre los grafos seleccionados, el resultado estructuralmente más complejo es \texttt{orbital\_pca2\_cubes10\_overlap0p35}. Este grafo tiene 124 nodos, 196 aristas y $\beta_1=96$, el valor más alto de ciclos entre los grafos seleccionados. Esta combinación lo convierte en el principal candidato para inspección científica, porque presenta una topología rica sin depender fuertemente de imputación.

La importancia del grafo orbital no debe entenderse como prueba definitiva de una estructura física real. Más bien, indica que el espacio definido por \texttt{pl\_orbper} y \texttt{pl\_orbsmax} produce un grafo con alta conectividad y múltiples ciclos. Esto sugiere que la arquitectura orbital del catálogo contiene regiones y transiciones interesantes, pero su interpretación final requiere revisar imputación, sesgo observacional y coherencia astrofísica.

\begin{table}[H]
\centering
\small
\begin{tabular}{p{0.43\textwidth}rrrr}
\toprule
\textbf{Grafo} & \textbf{Nodos} & \textbf{Aristas} & $\boldsymbol{\beta_1}$ & \textbf{Imputación media} \\
\midrule
\texttt{phys\_min\_pca2\_cubes10\_overlap0p35} & 66 & 105 & 54 & 0.017 \\
\texttt{phys\_density\_pca2\_cubes10\_overlap0p35} & 67 & 102 & 52 & 0.002 \\
\texttt{orbital\_pca2\_cubes10\_overlap0p35} & 124 & 196 & 96 & 0.011 \\
\texttt{joint\_no\_density\_pca2\_cubes10\_overlap0p35} & 62 & 134 & 81 & 0.176 \\
\texttt{joint\_pca2\_cubes10\_overlap0p35} & 56 & 126 & 77 & 0.152 \\
\texttt{thermal\_pca2\_cubes10\_overlap0p35} & 121 & 150 & 67 & 0.588 \\
\bottomrule
\end{tabular}
\caption{Métricas principales de los grafos seleccionados.}
\label{tab:selected_graph_metrics}
\end{table}

La Tabla~\ref{tab:selected_graph_metrics} muestra que la complejidad topológica por sí sola no basta para seleccionar un grafo como evidencia fuerte. El grafo térmico, por ejemplo, también presenta una estructura considerable, con 121 nodos, 150 aristas y $\beta_1=67$. Sin embargo, su imputación media nodal es 0.588, mucho mayor que la del grafo orbital. Esto significa que la estructura térmica puede ser informativa, pero no debe interpretarse con la misma confianza que una estructura sostenida por variables con baja imputación.

En contraste, el grafo orbital combina dos propiedades deseables: alta complejidad y baja imputación. Su imputación media nodal es 0.011, lo que lo vuelve más confiable que el grafo térmico para inspección científica inicial. Por esta razón, el espacio orbital se considera el resultado principal del análisis Mapper.

Los espacios físicos muestran una estructura más moderada. El grafo \texttt{phys\_min} tiene $\beta_1=54$, mientras que \texttt{phys\_density} tiene $\beta_1=52$. La diferencia es pequeña, pero relevante: al agregar densidad, la complejidad no aumenta. Esto sugiere que \texttt{pl\_dens} no introduce una nueva ramificación fuerte en el espacio físico, sino que parece regularizar o compactar ligeramente la estructura. Esta lectura es coherente con el hecho de que la densidad fue derivada en gran parte a partir de masa y radio, por lo que no debe tratarse como evidencia completamente independiente.

Una comparación similar aparece en los espacios conjuntos. El grafo \texttt{joint\_no\_density} tiene $\beta_1=81$ e imputación media de 0.176, mientras que \texttt{joint} tiene $\beta_1=77$ e imputación media de 0.152. Al agregar densidad, la complejidad conjunta disminuye ligeramente. Esto refuerza la lectura de que \texttt{pl\_dens} funciona más como una variable de control o regularización que como una fuente independiente de estructura topológica nueva.

\begin{table}[H]
\centering
\small
\begin{tabular}{p{0.32\textwidth}rrp{0.34\textwidth}}
\toprule
\textbf{Comparación} & \textbf{$\Delta \beta_1$} & \textbf{Cambio} & \textbf{Interpretación} \\
\midrule
\texttt{phys\_density} $-$ \texttt{phys\_min} & $52-54=-2$ & Disminuye & La densidad no aumenta la ramificación física. \\
\texttt{joint} $-$ \texttt{joint\_no\_density} & $77-81=-4$ & Disminuye & La densidad compacta ligeramente el espacio conjunto. \\
\bottomrule
\end{tabular}
\caption{Efecto de agregar \texttt{pl\_dens} en los espacios físico y conjunto.}
\label{tab:density_effect_mapper}
\end{table}

El resultado térmico debe leerse de forma separada. Aunque \texttt{thermal\_pca2\_cubes10\_overlap0p35} presenta una estructura no trivial, su alta imputación media nodal indica que gran parte de la forma observada puede depender de valores completados estadísticamente. Por ello, el espacio térmico se conserva como evidencia exploratoria, pero no como base principal para una conclusión física fuerte.

En conjunto, los resultados principales de Mapper sugieren tres lecturas. Primero, el espacio orbital es el candidato más informativo porque combina alta complejidad topológica y baja imputación. Segundo, los espacios físicos son útiles para interpretación, pero la densidad debe tratarse como control debido a su dependencia algebraica con masa y radio. Tercero, el espacio térmico muestra estructura, pero su alta dependencia de imputación obliga a una interpretación cautelosa. Estas observaciones preparan la siguiente etapa del análisis: interpretar los grafos seleccionados y auditar si sus regiones están asociadas con propiedades astrofísicas o con sesgos observacionales.